% Aula 03 --- de que tamanho é o viés da validação cruzada, e para que lado.
A validação cruzada estima o risco de um procedimento --- mas não exatamente o
risco que interessa. Há um descompasso embutido, e a boa notícia é que ele tem
sinal conhecido.
\begin{itens}
  \item \pts{1,0} Cada $g_{-j}$ é ajustado com quantas observações? Usando isso,
        argumente que
        $\E\big[\riscohat_{k\text{-CV}}\big]=\risco_{n-n/k}$, onde $\risco_m$
        denota o risco esperado do procedimento treinado com $m$ observações.

        A conta é mais curta do que parece: tome a esperança dentro da soma dupla
        e pergunte-se, para uma parcela qualquer, o que a observação $i$ representa
        para o ajuste $g_{-j}$ que a está prevendo.
  \item \pts{0,8} O alvo que interessa é $\risco_n$, o risco do procedimento
        treinado com a amostra inteira. Supondo que $\risco_m$ decresça com $m$ ---
        mais dados não pioram ---, conclua que a validação cruzada é
        \textbf{pessimista}.
  \item \pts{0,7} Ordene $k=2$, $k=10$ e $k=n$ quanto ao tamanho desse viés,
        justificando pelo número de observações de treino em cada caso.
  \item \pts{1,0} Se o LOOCV tem o menor viés, por que $k=5$ e $k=10$ são os
        valores usuais? Responda pensando em \emph{para que} a validação cruzada
        costuma ser usada.
\end{itens}

\begin{solucao}
\textbf{(a)} O ajuste $g_{-j}$ usa todas as observações menos as da dobra $L_j$,
isto é, $n-\abs{L_j}$ pontos. Com dobras de tamanho $n/k$, são $n-n/k$.

Tome a esperança na definição:
\[
  \E\big[\riscohat_{k\text{-CV}}\big]
   = \frac1n\sum_{j=1}^{k}\sum_{i\in L_j}
     \E\Big[\big(Y_i-g_{-j}(\X_i)\big)^2\Big].
\]
Agora olhe uma parcela isolada. A observação $i$ pertence a $L_j$, e $L_j$ foi
justamente a dobra retirada para treinar $g_{-j}$ --- então $i$ \textbf{não}
participou daquele ajuste. Para $g_{-j}$, ela é uma observação nova e
independente.

Mas "erro esperado de um ajuste numa observação nova" é exatamente a definição de
risco. Como $g_{-j}$ foi treinado com $n-n/k$ pontos, essa esperança vale
$\risco_{n-n/k}$ --- e vale o mesmo para toda parcela, por simetria, já que as
observações são i.i.d.

São $n$ parcelas no total (cada observação aparece em exatamente uma dobra), o
fator $1/n$ cancela, e sobra $\E[\riscohat_{k\text{-CV}}]=\risco_{n-n/k}$.

\smallskip
\textbf{(b)} Como $n-n/k<n$ e $\risco_m$ é decrescente em $m$, temos
$\risco_{n-n/k}\ge\risco_n$. Logo
\[
  \E\big[\riscohat_{k\text{-CV}}\big] = \risco_{n-n/k} \;\ge\; \risco_n .
\]
A validação cruzada superestima, em média, o risco do procedimento treinado com
tudo.

A origem do viés cabe numa frase: \textbf{a CV mede o desempenho de um modelo
treinado com menos dados do que aquele que você vai de fato usar.}

E o sinal conhecido é uma vantagem, não um defeito. Se a CV diz que o modelo é
bom, ele é pelo menos aquilo --- o erro está do lado seguro.

\smallskip
\textbf{(c)} O viés é tanto maior quanto menos observações sobram para treinar:

\begin{center}\small
\begin{tabular}{@{}lll@{}}
\toprule
 & observações de treino & viés \\
\midrule
$k=2$   & $n/2$    & o maior \\
$k=10$  & $0{,}9n$ & pequeno \\
$k=n$ (LOOCV) & $n-1$ & o menor \\
\bottomrule
\end{tabular}
\end{center}

Com $k=2$, metade dos dados é descartada em cada rodada, e o que se estima é o
risco de um procedimento com metade da amostra --- que pode ser bem pior. Com o
LOOCV, cada ajuste usa $n-1$ pontos, praticamente a amostra inteira, e o alvo
estimado é praticamente o certo.

\smallskip
\textbf{(d)} Porque o viés atrapalha menos do que parece para o uso mais comum, e
o LOOCV custa caro.

Se o objetivo for \textbf{reportar} o risco, o viés importa de verdade: o número
sai sistematicamente alto, e aí um $k$ grande ajuda.

Só que o uso mais frequente é \textbf{escolher} um hiperparâmetro --- descobrir
onde a curva $\theta\mapsto\riscohat(\theta)$ tem mínimo. Para isso o que importa
é a \emph{forma} da curva, não o nível dela. Um viés que desloca todos os valores
para cima de maneira parecida preserva a posição do mínimo, e a escolha não muda.

Some a isso o custo: o LOOCV exige $n$ ajustes e, fora do caso linear, não tem
atalho --- justamente o que o exercício do atalho mostra é que a economia só
existe quando o ajuste é linear. Com $k=10$ o viés já é pequeno e o custo é dez
ajustes, não mil.

Uma ressalva, para fechar: o argumento clássico diz também que $k$
grande aumenta a \emph{variância} da estimativa, porque as dobras ficam muito
parecidas entre si. É o que o deck da Aula~03 ensina, e é o que se costuma repetir
--- ainda que a medição desse efeito seja menos consensual do que o argumento
sugere.
\end{solucao}
